\section{一粒子の状態}
\label{sec:state}

地位は\dstrec{}である。型は \texttt{Algebra.Operations.TorsionParams} であり、場は次の二つである。

\begin{verbatim}
struct TorsionParams:
    alpha: array[3] of real   # usual rapidity phi
    beta:  array[3] of real   # dual angle theta
\end{verbatim}

添字 \(a\in\{0,1,2\}\) は Lean の \texttt{Fin 3} である。三軸は、書かれた作用の中では別々に積分する。軸を混ぜる項を、この記録のソルバへ足さない。

\subsection{共通ラピディティと不整合}

一軸ごとに
\[
\sigma_a=\alpha_a+\beta_a,\qquad
\delta_a=\alpha_a-\beta_a
\]
と置く。逆は
\[
\alpha_a=\frac{\sigma_a+\delta_a}{2},\qquad
\beta_a=\frac{\sigma_a-\delta_a}{2}
\]
である。\texttt{DualRotorFlow.usualFrom} と \texttt{dualFrom} は、時刻の関数 \(\sigma,\delta\) にこの一次式を点ごとに適用する。プログラムは \(\alpha,\beta\) と \(\sigma,\delta\) を別々の状態として持たない。\(\sigma(t)\) と \(\delta(t)\) を閉形式で計算し、角はこの一次式で作る。

\subsection{ねじれスカラーと符号なし質量}

実装する式は \texttt{Invariant.J} と \texttt{Invariant.mass} である。
\begin{align}
J(p)&=\frac12\sum_{a=0}^{2}\bigl(\alpha_a^2-\beta_a^2\bigr)
=\frac12\sum_{a=0}^{2}\delta_a\sigma_a,\\
M(p)&=\frac12\sum_{a=0}^{2}\bigl(\alpha_a^2+\beta_a^2\bigr)
=\frac14\sum_{a=0}^{2}\bigl(\sigma_a^2+\delta_a^2\bigr).
\end{align}
生成子のキリング形式を \(16\) で割った値は \(J\) ではない。恒等式は
\[
J(p)=4\cdot\Bigl(\frac1{16}\,\texttt{omegaTorsionGeneratorKilling}(p)\Bigr)
\]
であり、右辺の括弧は \(\frac18\sum_a(\alpha_a^2-\beta_a^2)\) である。係数 \(8\sum_a(\alpha_a^2-\beta_a^2)\) も \(J\) ではない。

正規化は
\[
J_{\mathrm{norm}}(p)=\frac{8}{3\pi^2}J(p)
\]
であり、識別子は \texttt{JNormalized} である。係数は正なので、\(|J_{\mathrm{norm}}|\le 1\) と \(|J|\le 3\pi^2/8\) はすべての状態で同値である。許容錐の上では、この不等式が成り立つ。定理は \texttt{torsion\_bound\_continuous} と \texttt{torsion\_bound\_raw\_continuous\_pi\_sq} である。錐の上で \(|J_{\mathrm{norm}}|=1\) となるのは、全軸が純双曲か全軸が純楕円かのときである。

\(J(p)=0\) は \(\sum_a\alpha_a^2=\sum_a\beta_a^2\) と同値である。各軸で \(\delta_a=0\) となることは、この等式の特別な場合である。軸をまたいで平方が打ち消し合っても \(J\) は消える。

一軸だけが非零のとき、\(J=\frac12\delta\sigma\)、\(M=\frac14(\sigma^2+\delta^2)\) はその軸への制限である。書かれた作用の剛性ポテンシャルは別の量で、
\[
\texttt{stiffnessPotential}(m,p)=\frac m2\sum_a(\alpha_a-\beta_a)^2
\]
である。これを \(J\) へ代入しない。一軸に限れば、\(m\neq 0\) のとき両者が \(m\) 倍で一致するのは \(\alpha=\beta\) または \(\beta=0\) のときに限る。

\(J=0\) は相対ローターが単位元であることを意味しない。一軸で \(\delta=0\) かつ \(\sigma\neq 0\) なら二つの角は等しく、\(\alpha=\beta=\sigma/2\neq 0\)、\(J=0\)、\(M=\sigma^2/4>0\) である。このとき相対ローターは単位元ではない。証人は \texttt{axisParams}\(1\,1\) であり、定理は \texttt{RelativeRotor.J\_zero\_not\_relativeRotor\_one} である。

\subsection{許容錐}

\texttt{Admissible.IsAdmissibleContinuous} は、各軸で
\[
\alpha_a\ge 0,\qquad \beta_a\ge 0,\qquad \alpha_a+\beta_a\le\frac\pi2
\]
が成り立つことである。\(\sigma,\delta\) へ訳すと、各軸で
\[
\sigma_a\in\Bigl[0,\frac\pi2\Bigr],\qquad |\delta_a|\le\sigma_a
\]
である。主枝の絶対値条件 \(|\alpha_a+\beta_a|\le\pi/2\) だけでは \(J\) は抑えられない。検査関数は連続許容の三つ組を見る。

許容錐を外れた状態の扱いは第\ref{sec:numerics}章で拒否とする。錐の中へ射影する関数は記録に無い。

\subsection{一つの質量にまとめない}

\texttt{RestMass} は、次の三つを一つの数にしない。

\begin{enumerate}
\item 剛性ポテンシャル \texttt{stiffnessPotential}。係数 \(m\) は入力である。式は前節のとおりである。
\item 符号なし質量 \(M\)。ラピディティのユークリッド二次形式である。
\item 自由落下。\(F=ma\) という識別は主張されていない。加速度の上限も主張されていない。
\end{enumerate}

段の更新に慣性パラメータが要るなら、呼び出し側が別に渡す。剛性 \(m\) でも \(M\) でもない。詳細は第\ref{sec:contract}章である。
